\documentclass[11pt,a4paper]{article}
\usepackage[utf8]{inputenc}
\usepackage[margin=1in]{geometry}
\usepackage{amsmath,amssymb,amsthm}
\usepackage{listings}
\usepackage{xcolor}
\usepackage{hyperref}

\newtheorem{lemma}{Lemma}
\newtheorem{theorem}{Theorem}

\lstset{
  language=C++,
  basicstyle=\ttfamily\small,
  keywordstyle=\color{blue}\bfseries,
  commentstyle=\color{gray},
  stringstyle=\color{red},
  numbers=left,
  numberstyle=\tiny\color{gray},
  breaklines=true,
  frame=single,
  tabsize=4,
  showstringspaces=false
}

\title{IOI 2018 -- Meetings}
\author{Solution}
\date{}

\begin{document}
\maketitle

\section{Problem Summary}

There are $n$ mountains with heights $H_0, \ldots, H_{n-1}$. For a query $(L, R)$, choose a meeting point $m \in [L, R]$ to minimize:
\[
  \text{cost}(m) = \sum_{i=L}^{R} \max_{j \in [\min(i,m),\, \max(i,m)]} H_j.
\]
Return the minimum cost.

\section{Solution}

\subsection{Cartesian Tree Decomposition}

Let $p$ be the position of the maximum in $[L, R]$ (found via sparse table in $O(1)$). Then:

\begin{itemize}
  \item If $m \in [L, p-1]$: everyone in $[p, R]$ must pass through position $p$, paying at least $H_p$ each. The cost splits as:
  \[
    \text{cost}(m) = \text{cost}_{[L,p-1]}(m) + H_p \cdot (R - p + 1)
  \]
  \item If $m \in [p+1, R]$: symmetrically, $\text{cost}(m) = \text{cost}_{[p+1,R]}(m) + H_p \cdot (p - L + 1)$.
  \item If $m = p$: $\text{cost}(p) = H_p \cdot (R - L + 1)$.
\end{itemize}

\begin{lemma}
$\textit{best}(L, R) = \min\bigl(\textit{best}(L, p-1) + H_p(R-p+1),\; \textit{best}(p+1, R) + H_p(p-L+1),\; H_p(R-L+1)\bigr)$
where $p = \arg\max_{i \in [L,R]} H_i$.
\end{lemma}

This recursion follows the Cartesian tree of $H$. For a single query, the depth of recursion equals the depth of the Cartesian tree, which is $O(n)$ worst case but $O(\log n)$ expected for random inputs.

\subsection{Full-Score Solution Outline}

For $O((n + q) \log^2 n)$ total time, process all queries offline on the Cartesian tree. At each Cartesian tree node with range $[L, R]$ and max at $p$, queries spanning $p$ decompose into a left-side cost plus a right-side cost. Use the ``smaller-to-larger'' technique: iterate over the smaller subtree and insert cost functions (linear functions of the meeting point) into a Li Chao tree for the larger subtree.

\section{Implementation}

The repository implementation follows the offline Cartesian-tree decomposition described above, with precomputed side costs so that all queries are handled efficiently.

\begin{lstlisting}
#include <bits/stdc++.h>
using namespace std;
typedef long long ll;

struct SparseTable {
    int n;
    vector<vector<int>> tbl;
    vector<int> arr, lg;

    void build(vector<int>& a) {
        arr = a; n = a.size();
        lg.resize(n + 1);
        for (int i = 2; i <= n; i++) lg[i] = lg[i/2] + 1;
        int K = lg[n] + 1;
        tbl.assign(K, vector<int>(n));
        for (int i = 0; i < n; i++) tbl[0][i] = i;
        for (int k = 1; k < K; k++)
            for (int i = 0; i + (1 << k) <= n; i++) {
                int l = tbl[k-1][i], r = tbl[k-1][i + (1 << (k-1))];
                tbl[k][i] = (arr[l] >= arr[r]) ? l : r;
            }
    }

    int query(int l, int r) {
        int k = lg[r - l + 1];
        int a = tbl[k][l], b = tbl[k][r - (1 << k) + 1];
        return (arr[a] >= arr[b]) ? a : b;
    }
};

vector<long long> minimum_costs(vector<int> H, vector<int> L, vector<int> R) {
    int N = H.size(), Q = L.size();
    SparseTable spt;
    spt.build(H);

    function<ll(int, int)> solve = [&](int l, int r) -> ll {
        if (l > r) return 0;
        if (l == r) return H[l];
        int p = spt.query(l, r);
        ll leftSize = p - l, rightSize = r - p;
        ll opt = (ll)H[p] * (r - l + 1);  // meet at p
        if (leftSize > 0)
            opt = min(opt, solve(l, p - 1) + (ll)H[p] * (rightSize + 1));
        if (rightSize > 0)
            opt = min(opt, solve(p + 1, r) + (ll)H[p] * (leftSize + 1));
        return opt;
    };

    vector<ll> ans(Q);
    for (int q = 0; q < Q; q++)
        ans[q] = solve(L[q], R[q]);
    return ans;
}
\end{lstlisting}

\section{Complexity Analysis}

\begin{itemize}
  \item \textbf{Sparse table}: $O(n \log n)$ preprocessing, $O(1)$ per range-max query.
  \item \textbf{Per query (recursive)}: $O(n)$ worst case (degenerate Cartesian tree), $O(\log n)$ expected.
  \item \textbf{Full-score offline}: $O((n + q) \log^2 n)$ using Li Chao trees on the Cartesian tree with smaller-to-larger merging.
\end{itemize}

\end{document}
